% Pandoc template for the fed_fraud documents.
% Based on proposal_gen/templates/detailed_proposal_template/detailed_1_classic.tex
% (Classic report): centered title page, table of contents, ruled section
% headings and the Geodesic house footer. Fonts are Latin Modern (lmodern).
% Two layouts: the full report (title page and contents) and, with
% -V brief, a compact one-page-brief header in the same style.
\documentclass[letterpaper,11pt]{article}
\usepackage[T1]{fontenc}
\usepackage[utf8]{inputenc}
\usepackage{lmodern}
\usepackage[letterpaper,top=0.9in,bottom=1.2in,left=1in,right=1in,footskip=0.6in]{geometry}
\usepackage[table]{xcolor}
\usepackage{geodesicbrand}
\usepackage{amsmath,amssymb}
\usepackage{enumitem}
\usepackage{array}
\usepackage{calc}
\usepackage{longtable}
\usepackage{makecell}
\usepackage[most]{tcolorbox}
\usepackage{titlesec}
\usepackage{needspace}
\usepackage{etoolbox}
\usepackage{hyperref}
\hypersetup{colorlinks=true,urlcolor=coolblack,linkcolor=coolblack,filecolor=geoblue,
  pdftitle={Policy and legal track},pdfauthor={\geofirm}}
\urlstyle{same}

\setlist{itemsep=2pt,leftmargin=1.4em}
\setlength{\parindent}{0pt}
\setlength{\parskip}{6pt}
\providecommand{\tightlist}{\setlength{\itemsep}{0pt}\setlength{\parskip}{0pt}}
\setcounter{secnumdepth}{0}

% Section style: house rule under each heading
\titleformat{\section}{\Large\bfseries\color{coolblack}}{\thesection}{0.6em}{}[{\color{coolblack}\titlerule[0.5pt]}]
\titleformat{\subsection}{\large\bfseries\color{geomid}}{\thesubsection}{0.6em}{}
\titleformat{\subsubsection}{\normalsize\bfseries\color{geomid}}{\thesubsubsection}{0.6em}{}
\titlespacing*{\section}{0pt}{10pt}{6pt}

% Tables are written by latex/filters.lua (navy header, banded rows)
\setlength{\LTpre}{6pt}\setlength{\LTpost}{6pt}

\begin{document}

% BRIEF HEADER
\thispagestyle{geodesic}
\begin{minipage}[c]{0.13\textwidth}\includegraphics[width=\linewidth]{icon.png}\end{minipage}\hfill
\begin{minipage}[c]{0.84\textwidth}
  {\textsc{\large \geofirm}}\par\vspace{2pt}
  {\small\color{geogray}Federal fraud expansion\enspace\textperiodcentered\enspace When: \compiletime}
\end{minipage}\par\vspace{6pt}
{\color{coolblack}\rule{\linewidth}{0.3mm}}\par\vspace{2pt}
{\LARGE\bfseries\color{coolblack} Policy and legal track}\par
\vspace{2pt}{\large\color{geomid} What the law allows, and what would unlock more}\par
\vspace{-2pt}{\color{coolblack}\rule{\linewidth}{0.3mm}}\par\vspace{4pt}
\pagestyle{geodesic}

\textbf{Status:} Draft v0.1 (2026-11-04), for Steve's review. Extension
5 of 10 (see the \href{../extensions.pdf}{extensions overview}).
\textbf{Not legal advice.} This is a working map for planning; counsel
must review it before any agreement is signed, and every citation should
be checked against the current U.S. Code and OMB guidance.

\hypertarget{bottom-line}{%
\section*{Bottom line}\label{bottom-line}}
\addcontentsline{toc}{section}{Bottom line}

Most of what this project needs is already lawful; it is slow, not
forbidden. Public data needs no permission. Inspectors General can see
their agency's records today. Treasury's Do Not Pay can match payments
against death and exclusion data under the Payment Integrity Information
Act. Cross-agency matching of people is allowed with a computer matching
agreement, which takes months but is routine. The binding limits are
narrower: tax data, the National Directory of New Hires, state-held
benefit data, and the time it takes to sign agreements. Five legislative
or policy changes would remove most of the delay; they are in section 6.

\hypertarget{the-laws-that-matter}{%
\section*{1. The laws that matter}\label{the-laws-that-matter}}
\addcontentsline{toc}{section}{1. The laws that matter}

\begingroup\small\setlength{\tabcolsep}{4pt}\renewcommand{\arraystretch}{1.25}
\rowcolors{2}{white}{geolight}
\begin{longtable}[l]{@{}>{\raggedright\arraybackslash}p{\dimexpr 0.333\linewidth-2\tabcolsep}>{\raggedright\arraybackslash}p{\dimexpr 0.333\linewidth-2\tabcolsep}>{\raggedright\arraybackslash}p{\dimexpr 0.333\linewidth-2\tabcolsep}@{}}
\rowcolor{coolblack}\leavevmode\color{white}\bfseries Law or rule & \leavevmode\color{white}\bfseries What it does for this project & \leavevmode\color{white}\bfseries Constraint \\
\endhead
Payment Integrity Information Act of 2019 (PIIA), 31 U.S.C. 3351-3358 & Requires agencies to estimate and reduce improper payments; establishes the Do Not Pay initiative (3354) and the fraud risk and data analytics duties that replaced the Fraud Reduction and Data Analytics Act of 2015 (3357) & Do Not Pay matching is limited to the databases and uses the statute and OMB allow \\
OMB M-21-19 (Appendix C to Circular A-123) & How agencies measure improper payments and use Do Not Pay & Guidance, updated periodically; check for a newer memo \\
Privacy Act of 1974, 5 U.S.C. 552a & Governs federal records about individuals: systems of records notices, routine uses, disclosure limits & Applies to individuals, not businesses; sole proprietors are a grey area \\
Computer Matching and Privacy Protection Act of 1988 (amending the Privacy Act, 552a(o)-(u)) & Allows matching of federal benefit and payroll records between agencies under a written matching agreement approved by each agency's Data Integrity Board & Agreements run 18 months plus a 12-month renewal; requires independent verification and notice before adverse action. PIIA lets Do Not Pay agreements run up to three years \\
Inspector General Act of 1978, 5 U.S.C. 401-424 & IGs have access to all agency records and can subpoena others & Investigative matches by an IG are excepted from the matching rules \\
E-Government Act of 2002, section 208 & Privacy impact assessments for new systems that hold personal data & Needed before any internal-data system goes live \\
26 U.S.C. 6103 & Tax return confidentiality & Tax data can leave the IRS only in listed cases; work must sit inside IRS or TIGTA \\
42 U.S.C. 653 (National Directory of New Hires) & Employment and wage data used for child support & Access limited to listed programs; Do Not Pay has no general access \\
Social Security Act 205(r) and the Consolidated Appropriations Act, 2021 & Gave Do Not Pay access to SSA's full death file for three years from December 2023 & The access was time-limited; check whether it has been extended \\
False Claims Act, 31 U.S.C. 3729-3733 & Civil liability (treble damages and penalties) for false claims; private relators can file qui tam suits & Lead packages should be written so DOJ can use them \\
Administrative False Claims Act (FY2025 NDAA, replacing the Program Fraud Civil Remedies Act) & Lets agencies pursue false claims up to \$1M administratively & Agency must have set up the process; check implementation \\
Whistleblower protections, 5 U.S.C. 2302(b)(8) and 41 U.S.C. 4712 & Protect federal employees and contractor or grantee employees who report fraud & Tip triage must flag these tipsters (see \href{../playbook/tip_triage.pdf}{tip triage}) \\
Paperwork Reduction Act, 44 U.S.C. 3501 & Information collection from the public needs OMB approval & Applies if a pilot collects new data from the public; using existing records does not \\
Executive Order 14243 (March 2025), ``Stopping Waste, Fraud, and Abuse by Eliminating Information Silos'' & Directs agencies to give designated officials full access to unclassified records & Does not override statutes such as 6103 or the Privacy Act \\
Executive Order 14249 (March 2025), ``Protecting America's Bank Account Against Fraud, Waste, and Abuse'' & Directs Treasury to verify payments before disbursement, including through Do Not Pay & Implementation details in Treasury guidance \\
2026 anti-fraud task force order (the america.gov/fraud initiative) & Makes cross-agency fraud work a stated priority & Confirm the order's number and text before quoting it \\
\end{longtable}\endgroup

\hypertarget{what-we-can-do-today-by-data-type}{%
\section*{2. What we can do today, by data
type}\label{what-we-can-do-today-by-data-type}}
\addcontentsline{toc}{section}{2. What we can do today, by data type}

\begingroup\small\setlength{\tabcolsep}{4pt}\renewcommand{\arraystretch}{1.25}
\rowcolors{2}{white}{geolight}
\begin{longtable}[l]{@{}>{\raggedright\arraybackslash}p{\dimexpr 0.308\linewidth-2\tabcolsep}>{\raggedright\arraybackslash}p{\dimexpr 0.308\linewidth-2\tabcolsep}>{\raggedright\arraybackslash}p{\dimexpr 0.308\linewidth-2\tabcolsep}>{\raggedright\arraybackslash}p{\dimexpr 0.077\linewidth-2\tabcolsep}@{}}
\rowcolor{coolblack}\leavevmode\color{white}\bfseries Data & \leavevmode\color{white}\bfseries Route & \leavevmode\color{white}\bfseries Paperwork & \leavevmode\color{white}\bfseries Time \\
\endhead
Public data (USAspending, SAM.gov, CMS files, SBA loan files, exclusion lists) & None needed & Terms of use & Now \\
Agency records about businesses & Program office data use agreement or OIG partnership & Data use agreement, security review & 2-6 months \\
Agency records about individuals, one agency & OIG partnership (IG Act) or routine use in the agency's system of records notice & Memorandum of understanding, privacy review & 3-6 months \\
Matching individuals across agencies & Computer matching agreement, or Do Not Pay, or an IG investigative match & Agreement, Data Integrity Board approval, Federal Register notice & 6-12 months \\
State-run program data (SNAP, UI, Medicaid) & State data sharing agreement; NASWA Integrity Data Hub for UI; T-MSIS through CMS for Medicaid & State agreement per state & 3-9 months each \\
Tax data & Only inside IRS or TIGTA & --- & --- \\
\end{longtable}\endgroup

\hypertarget{the-computer-matching-path-step-by-step}{%
\section*{3. The computer matching path, step by
step}\label{the-computer-matching-path-step-by-step}}
\addcontentsline{toc}{section}{3. The computer matching path, step by
step}

\begin{enumerate}
\def\labelenumi{\arabic{enumi}.}
\tightlist
\item
  Identify the systems of records on both sides and confirm the match is
  a ``matching program'' (federal benefit or payroll records, automated
  comparison).
\item
  Draft the matching agreement: purpose, legal authority, justification
  and expected results (a cost-benefit analysis, which the
  \href{../proposals/roi_model.pdf}{ROI model} supplies), records and
  fields, notice procedures, verification and due process, retention,
  security.
\item
  Data Integrity Board approval at each agency.
\item
  Report to OMB and Congress; publish the notice in the Federal Register
  at least 30 days before matching starts.
\item
  Independent verification of every hit, and notice with at least 30
  days to contest, before any adverse action.
\end{enumerate}

Ways to shorten it: run the first match as an IG investigative match;
route through Do Not Pay, which already holds agreements and data; reuse
an existing agreement's terms.

\hypertarget{privacy-commitments-we-make-in-every-proposal}{%
\section*{4. Privacy commitments we make in every
proposal}\label{privacy-commitments-we-make-in-every-proposal}}
\addcontentsline{toc}{section}{4. Privacy commitments we make in every
proposal}

\begin{itemize}
\tightlist
\item
  Minimum necessary fields; identifiers hashed across boundaries.
\item
  Analysis inside the data owner's authorized environment.
\item
  No adverse action on a model score; a person decides, with
  verification and notice.
\item
  Model documentation and bias testing for every score.
\item
  Publication only in aggregate, with small cells suppressed.
\end{itemize}

\hypertarget{risks-to-manage}{%
\section*{5. Risks to manage}\label{risks-to-manage}}
\addcontentsline{toc}{section}{5. Risks to manage}

\begin{itemize}
\tightlist
\item
  \textbf{Sole proprietors:} business records that are really records
  about individuals bring in the Privacy Act. Treat them as personal
  data.
\item
  \textbf{Mission creep:} data shared for fraud detection must not be
  reused for other purposes without new authority.
\item
  \textbf{Model error and disparate impact:} a screen that flags one
  community more than its real fraud rate warrants harms people and the
  case. Bias checks are part of release, not an afterthought.
\item
  \textbf{Litigation holds and discovery:} leads may be discoverable;
  write them as factual reasons, not conclusions.
\end{itemize}

\hypertarget{legislative-and-policy-changes-that-would-unlock-better-data}{%
\section*{6. Legislative and policy changes that would unlock better
data}\label{legislative-and-policy-changes-that-would-unlock-better-data}}
\addcontentsline{toc}{section}{6. Legislative and policy changes that
would unlock better data}

Ranked by value for effort.

\begin{enumerate}
\def\labelenumi{\arabic{enumi}.}
\tightlist
\item
  \textbf{Make Do Not Pay's access to SSA's full death file permanent}
  and extend it to state-administered programs. Single biggest win for
  payments to dead people (roadmap X-1).
\item
  \textbf{Give Do Not Pay access to the National Directory of New Hires}
  for benefit programs with income tests (UI, SNAP, SSI, Medicaid). Long
  proposed by Treasury and in past budget requests.
\item
  \textbf{Standing cross-agency analytics authority} like the Pandemic
  Response Accountability Committee's analytics center, but permanent
  and covering all programs, with the matching rules adapted (longer
  agreements, a single Data Integrity Board process).
\item
  \textbf{Require states to share eligibility and payment data for
  federally funded programs} with the federal program office and OIG in
  a standard format (T-MSIS-style for SNAP and UI).
\item
  \textbf{Payee bank account verification for every federal payment}
  before first disbursement (account-owner match), building on EO 14249.
\end{enumerate}

\hypertarget{questions-for-steve}{%
\section*{7. Questions for Steve}\label{questions-for-steve}}
\addcontentsline{toc}{section}{7. Questions for Steve}

\begin{enumerate}
\def\labelenumi{\arabic{enumi}.}
\tightlist
\item
  Do we have counsel who has negotiated a computer matching agreement?
\item
  Should we write a short white paper on changes 1 and 2 for the task
  force now, since both have prior support?
\end{enumerate}

\end{document}
